% ID: MAT-ALG-PER-007
% Titolo: Sconti successivi
% Disciplina: Matematica
% Area: Algebra
% Argomento: Percentuali
% Classe: Prima secondaria di secondo grado
% Difficolta: 3
% Tipo: laboratorio
% Tag: matematica, algebra, percentuali
% Fonte: laboratorio_percentuali_anonimizzato
% Provenienza: materiale fornito dal docente; testo anonimizzato; soluzione redatta per Exergo

\begin{esercizio}
Un articolo costa 100 euro e subisce uno sconto del $20\%$ seguito da uno del $10\%$. Calcola prezzo finale e sconto complessivo. Spiega inoltre perché due sconti successivi del $20\%$ non equivalgono al $40\%$.
\end{esercizio}

\begin{soluzione}
Il primo prezzo è 80 euro, il secondo $80\cdot0{,}90=72$ euro. Sconto complessivo $28\%$. Due sconti del $20\%$ danno $100\cdot0{,}8^2=64$, dunque $36\%$: il secondo sconto si applica a una base ridotta.
\end{soluzione}
